\documentclass[11pt,a4paper]{article}

% ---------- Paquetes ----------
\usepackage[utf8]{inputenc}
\usepackage[T1]{fontenc}
\usepackage[english]{babel}
\usepackage{amsmath,amssymb,amsthm}
\usepackage{mathtools}
\usepackage{geometry}
\usepackage{booktabs}
\usepackage{array}
\usepackage{enumitem}
\usepackage{hyperref}
\usepackage{xcolor}
\usepackage{listings}
\usepackage{algorithm}
\usepackage{algpseudocode}
\usepackage{caption}
\usepackage{subcaption}
\usepackage{fancyhdr}
\usepackage{titlesec}

\geometry{margin=1in}
\hypersetup{
    colorlinks=true,
    linkcolor=blue!50!black,
    citecolor=blue!50!black,
    urlcolor=blue!50!black,
    pdftitle={Entropy Arena: A Mathematical Framework for Comparing Randomness Sources},
    pdfauthor={Entropy Arena Project}
}

% ---------- Entornos ----------
\theoremstyle{plain}
\newtheorem{theorem}{Theorem}[section]
\newtheorem{lemma}[theorem]{Lemma}
\newtheorem{proposition}[theorem]{Proposition}
\newtheorem{corollary}[theorem]{Corollary}

\theoremstyle{definition}
\newtheorem{definition}[theorem]{Definition}
\newtheorem{example}[theorem]{Example}
\newtheorem{remark}[theorem]{Remark}

% ---------- Encabezado ----------
\pagestyle{fancy}
\fancyhf{}
\rhead{\small Entropy Arena}
\lhead{\small\thesection\ \textit{\nouppercase{\leftmark}}}
\cfoot{\thepage}

% ---------- Título ----------
\title{
    \vspace{-2cm}
    \rule{\linewidth}{0.4pt}\\[0.4cm]
    \textbf{\Large Entropy Arena}\\[0.3cm]
    \large A Mathematical Framework for the Comparative\\ 
    Analysis of Randomness Sources in Cryptographic Key Generation\\[0.4cm]
    \rule{\linewidth}{0.4pt}
}
\author{\normalsize Entropy Arena Project}
\date{\today}

\begin{document}
\maketitle

\begin{abstract}
\noindent
This document formalizes the mathematical foundations of \emph{Entropy Arena}, a
framework for the empirical and theoretical comparison of randomness sources used
in the generation of seeds and private keys for a multi-signature wallet. We
present rigorous definitions of the entropy measures employed (Shannon,
min-entropy, collision, and compression-based), the dynamical systems underlying
the mathematical generators (Brownian motion, logistic map, Lorenz attractor,
and cellular automata), the statistical battery used to evaluate the outputs
(frequency, runs, chi-squared, autocorrelation, and a subset of NIST SP~800-22),
the composite scoring functional, and the cryptographic pipeline that maps raw
entropy to BIP32 keys and finally to a $2$-of-$3$ multisignature script.
\end{abstract}

\tableofcontents
\newpage

% ============================================================
\section{Introduction}
% ============================================================

The security of any cryptographic system ultimately rests on the quality of the
randomness used to generate its keys. In the context of a multi-signature wallet,
this requirement is amplified: multiple independent entropy sources must be
combined, each ideally derived from a distinct physical or computational origin.
\emph{Entropy Arena} provides a common formal ground on which heterogeneous
generators --- computational, mathematical, and physically manual --- can be
compared under identical statistical scrutiny.

Let $\mathcal{M} = \{M_1, M_2, \dots, M_k\}$ be a finite set of entropy methods.
Each method $M_i$ is modeled as a stochastic map
\begin{equation}
    M_i : \mathbb{N} \to \{0,1\}^{8n}, \qquad M_i(n) = D_i,
\end{equation}
that produces a byte string $D_i$ of length $n$. The arena's objective is to
assign to each $M_i$ a score $S(M_i) \in [0,100]$ that reflects both the
statistical quality of the output and the \emph{honesty} of the entropy claimed.

% ============================================================
\section{Preliminaries: Measures of Entropy}
% ============================================================

\subsection{Shannon entropy}

\begin{definition}[Shannon entropy]
Let $X$ be a discrete random variable with alphabet $\mathcal{X}$ and probability
mass function $p(x) = \Pr[X = x]$. The Shannon entropy of $X$, in bits, is
\begin{equation}
    H_1(X) \;=\; -\sum_{x \in \mathcal{X}} p(x) \log_2 p(x).
    \label{eq:shannon}
\end{equation}
For a byte-valued variable, $\mathcal{X} = \{0, \dots, 255\}$ and
$H_1(X) \le 8$.
\end{definition}

Given an empirical byte string $D = (d_1, \dots, d_n)$, we estimate the
distribution through frequency counts $c_v = |\{i : d_i = v\}|$ for each
$v \in \{0,\dots,255\}$, giving $\hat{p}(v) = c_v / n$ and
\begin{equation}
    \hat{H}_1(D) \;=\; -\sum_{v=0}^{255} \hat{p}(v) \log_2 \hat{p}(v).
\end{equation}

\subsection{Min-entropy}

\begin{definition}[Min-entropy]
The min-entropy of $X$ is
\begin{equation}
    H_\infty(X) \;=\; -\log_2 \max_{x \in \mathcal{X}} p(x).
    \label{eq:minentropy}
\end{equation}
It lower-bounds the output of any extractor: even the most likely outcome
carries at most $H_\infty(X)$ bits of unpredictability.
\end{definition}

Empirically, $\hat{H}_\infty(D) = -\log_2 \max_v \hat{p}(v)$.

\begin{remark}
Min-entropy is the \emph{honest} measure for cryptographic purposes. A
cryptographic hash applied to any input produces an output whose Shannon entropy
approaches $8$ bits/byte, yet whose min-entropy is bounded by the entropy of the
input. Thus, $H_1$ alone is insufficient to judge a generator.
\end{remark}

\subsection{Collision entropy}

\begin{definition}[Collision entropy]
Also known as Rényi entropy of order $2$:
\begin{equation}
    H_2(X) \;=\; -\log_2 \sum_{x \in \mathcal{X}} p(x)^2.
    \label{eq:collision}
\end{equation}
It relates to the collision probability $P_c(X) = \sum_x p(x)^2$ via
$H_2(X) = -\log_2 P_c(X)$.
\end{definition}

For a uniform distribution over 256 values, $H_1 = H_2 = H_\infty = 8$ bits.

\subsection{Compression-based redundancy}

\begin{definition}[Compression ratio]
Let $\mathcal{C}$ denote a lossless compressor (here, DEFLATE). The compression
ratio of $D$ is
\begin{equation}
    \rho(D) \;=\; \frac{|\mathcal{C}(D)|}{|D|}.
    \label{eq:compress}
\end{equation}
Incompressible strings satisfy $\rho \approx 1$; highly redundant strings yield
$\rho \ll 1$.
\end{definition}

% ============================================================
\section{Entropy Sources}
% ============================================================

Each method $M_i$ belongs to one of four categories:
\begin{itemize}[leftmargin=1.4em]
    \item \textbf{Computational:} relies on the OS CSPRNG or on a non-secure PRNG.
    \item \textbf{Mathematical:} uses a deterministic dynamical system, seeded
          from the OS.
    \item \textbf{Physical--manual:} entropy gathered by the user without any
          digital intervention.
    \item \textbf{Hybrid:} manual entropy passed through a whitening or
          bias-correcting extractor.
\end{itemize}

\subsection{Computational methods}

\paragraph{CSPRNG (\texttt{secrets}, \texttt{os.urandom}).} Let
$G : \{0,1\}^\lambda \to \{0,1\}^{n}$ be a cryptographic pseudorandom generator.
Assuming the seed is uniform over $\{0,1\}^\lambda$ and $G$ is a secure PRG,
the output is computationally indistinguishable from uniform:
\begin{equation}
    \big|\Pr[\mathcal{A}(G(U_\lambda)) = 1] - \Pr[\mathcal{A}(U_n) = 1]\big|
    \;\le\; \varepsilon(\lambda)
\end{equation}
for every probabilistic polynomial-time adversary $\mathcal{A}$. The min-entropy
of the output, conditioned on the seed being secret, is $H_\infty = n$ bits.

\paragraph{Mersenne Twister.} The MT19937 generator maintains a state
$\mathbf{x} = (x_0, \dots, x_{623}) \in (\mathbb{Z}/2^{32}\mathbb{Z})^{624}$
and evolves via
\begin{equation}
    x_{k+n} \;=\; x_{k+m} \;\oplus\; \big((x_k^{u} \mid x_{k+1}^{l})\,A\big),
    \label{eq:mt}
\end{equation}
with $n=624$, $m=397$, and a constant twist matrix $A$. Its period is
$2^{19937}-1$, but its state is fully recoverable from $624$ consecutive
32-bit outputs. Consequently, if the seed is known,
\begin{equation}
    H_\infty(M_{\text{MT}}(n) \mid \text{seed}) \;=\; 0.
\end{equation}

\paragraph{UUIDv4.} A version-4 UUID embeds $122$ random bits out of $128$;
the remaining $6$ bits encode the version and variant. Hence
\begin{equation}
    H_\infty(\text{UUIDv4}) \;=\; 122 \text{ bits}, \qquad
    \frac{122}{16} \;=\; 7.625 \text{ bits/byte}.
\end{equation}

\subsection{Mathematical methods: deterministic dynamical systems}

All methods in this category are deterministic given an initial condition, and
their unpredictability derives entirely from the entropy injected into the seed.
We formalize this via the following invariant.

\begin{proposition}[Entropy conservation]
Let $f : \mathcal{S} \to \mathcal{S}$ be a deterministic map and let
$s_0 \sim \mu$ be a random initial condition. The output trajectory
$(s_0, f(s_0), f^2(s_0), \dots)$ has entropy bounded above by $H(s_0)$.
Formally, for any function $\phi$ of the trajectory,
\begin{equation}
    H(\phi(s_0, s_1, \dots)) \;\le\; H(s_0).
\end{equation}
\end{proposition}

\begin{proof}
The map $s_0 \mapsto (s_0, s_1, \dots)$ is injective on $\mathcal{S}$; applying
the data-processing inequality yields the result.
\end{proof}

\subsubsection{Brownian motion}

The standard Wiener process $\{W_t\}_{t \ge 0}$ satisfies $W_0 = 0$ and, for
$0 \le s < t$, $W_t - W_s \sim \mathcal{N}(0, t-s)$ with independent increments.
We discretize with step $\Delta t$ and increments
\begin{equation}
    \xi_k \;=\; \sqrt{\Delta t}\;\mathcal{N}(0,1), \qquad
    W_k \;=\; \sum_{j=1}^{k} \xi_j.
    \label{eq:bm}
\end{equation}
The bit stream is extracted from the sign of the increments:
\begin{equation}
    b_k \;=\; \mathbf{1}\!\left[\xi_k > 0\right].
    \label{eq:bm_bits}
\end{equation}
Since $\Pr[\xi_k > 0] = 1/2$, each bit carries at most $1$ bit of Shannon
entropy but, being a deterministic function of the seed, contributes
$H_\infty = 0$ \emph{new} entropy.

\subsubsection{Logistic map}

The logistic map on the unit interval is
\begin{equation}
    x_{k+1} \;=\; r\, x_k\,(1 - x_k), \qquad x_k \in (0,1), \quad r \in (0,4].
    \label{eq:logistic}
\end{equation}
For $r \in [r_\infty, 4]$ with $r_\infty \approx 3.56995$, the dynamics are
chaotic. The Lyapunov exponent is
\begin{equation}
    \lambda \;=\; \lim_{N \to \infty} \frac{1}{N} \sum_{k=0}^{N-1}
    \ln\big| r (1 - 2x_k) \big|.
    \label{eq:lyap}
\end{equation}
For $r = 4$, $\lambda = \ln 2$, indicating exponential divergence of nearby
trajectories. A burn-in of $K$ iterations ($K \approx 10^3$) discards the
transient. Bits are extracted by
\begin{equation}
    b_k \;=\; \mathbf{1}\!\left[x_k > \tfrac{1}{2}\right].
    \label{eq:logistic_bits}
\end{equation}

\subsubsection{Lorenz attractor}

The Lorenz system is the three-dimensional ODE
\begin{equation}
    \dot{x} = \sigma(y - x), \qquad
    \dot{y} = x(\rho - z) - y, \qquad
    \dot{z} = xy - \beta z,
    \label{eq:lorenz}
\end{equation}
with canonical parameters $\sigma = 10$, $\rho = 28$, $\beta = 8/3$. Numerical
integration uses explicit Euler with step $\Delta t = 10^{-2}$. Bits are
extracted from sign changes of the $x$-coordinate:
\begin{equation}
    b_k \;=\; \mathbf{1}\!\left[x_{k+1} > x_k\right].
    \label{eq:lorenz_bits}
\end{equation}
The system's largest Lyapunov exponent is $\lambda_1 \approx 0.906$
(natural log), confirming sensitive dependence on initial conditions.

\subsubsection{Rule 30 cellular automaton}

Let $c_i^{(t)} \in \{0,1\}$ denote the state of cell $i$ at time $t$ on a ring
of width $L$. The elementary cellular automaton with Wolfram rule $30$ updates
via
\begin{equation}
    c_i^{(t+1)} \;=\; c_{i-1}^{(t)} \;\oplus\;
    \big(c_i^{(t)} \;\lor\; c_{i+1}^{(t)}\big),
    \label{eq:rule30}
\end{equation}
where indices are taken modulo $L$. After $K$ warm-up steps, the first byte of
each row is extracted:
\begin{equation}
    b^{(t)} \;=\; \sum_{i=0}^{7} c_i^{(t)}\, 2^{7-i}.
\end{equation}
Rule 30 is known to be computationally universal and passes the majority of
Diehard tests, yet remains fully deterministic given the initial configuration.

\subsection{Physical--manual methods}

\subsubsection{Dice rolls}

A fair six-sided die produces outcomes uniformly distributed over
$\{1,\dots,6\}$, contributing
\begin{equation}
    h_{\text{die}} \;=\; \log_2 6 \;\approx\; 2.5850 \text{ bits per roll}.
    \label{eq:die_bits}
\end{equation}
For $n$ rolls with outcomes $(r_1, \dots, r_n)$, we form the integer
\begin{equation}
    N \;=\; \sum_{i=1}^{n} (r_i - 1)\, 6^{\,n-i},
    \label{eq:die_int}
\end{equation}
so that $N \in \{0, 1, \dots, 6^n - 1\}$. The total entropy is
\begin{equation}
    H_\infty(\text{dice}_n) \;=\; n \log_2 6 \;\approx\; 2.585\, n \text{ bits}.
    \label{eq:die_total}
\end{equation}
A cryptographic extractor $E : \mathbb{Z} \to \{0,1\}^{8m}$ (here SHA-256
applied iteratively) maps $N$ to $m$ bytes:
\begin{equation}
    D \;=\; E(N) \;=\; \mathrm{SHA256}(N) \,\|\, \mathrm{SHA256}(E_1) \,\|\, \cdots
\end{equation}
where $E_1 = \mathrm{SHA256}(N)$ and subsequent blocks are chained. The
extractor preserves min-entropy: $H_\infty(D) = \min(8m, n \log_2 6)$.

\subsubsection{Coin flips}

A fair coin yields $h_{\text{flip}} = 1$ bit per toss. For $n$ flips with
outcomes $f_i \in \{H, T\}$ we set
\begin{equation}
    b_i \;=\; \mathbf{1}[f_i = H],
\end{equation}
and pack the bits into bytes via
\begin{equation}
    d_j \;=\; \sum_{i=0}^{7} b_{8j + i}\, 2^{7-i}.
\end{equation}
The total min-entropy is $n$ bits.

\subsubsection{Bias correction: von Neumann extractor}

If the coin is biased with $\Pr[H] = p \ne 1/2$, the raw bits carry less than
$1$ bit each. The von Neumann extractor operates on non-overlapping pairs:
\begin{equation}
    (H,T) \mapsto 1, \qquad (T,H) \mapsto 0, \qquad (H,H), (T,T) \mapsto \varnothing.
    \label{eq:vonneumann}
\end{equation}
Each surviving bit is exactly fair. The acceptance probability per pair is
\begin{equation}
    q \;=\; 2 p (1 - p),
    \label{eq:vn_eff}
\end{equation}
so on average $n q / 2$ output bits are obtained from $n$ flips. The expected
number of flips required per output bit is
\begin{equation}
    \mathbb{E}[\text{flips/bit}] \;=\; \frac{1}{p(1-p)} \;\ge\; 4.
\end{equation}
For $p = 1/2$, this equals $4$; efficiency drops sharply as $|p - 1/2|$ grows.

\subsubsection{Deck shuffle}

A uniformly random permutation of $52$ cards has entropy
\begin{equation}
    H_\infty(\text{deck}) \;=\; \log_2(52!) \;\approx\; 225.581 \text{ bits},
    \label{eq:deck_bits}
\end{equation}
or approximately $4.338$ bits per card. Given a permutation
$\pi = (\pi_1, \dots, \pi_{52})$, we encode it via the Lehmer code:
\begin{equation}
    N \;=\; \sum_{i=1}^{52} c_i\, (52-i)!,
    \label{eq:lehmer}
\end{equation}
where
\begin{equation}
    c_i \;=\; \big|\{ j > i : \pi_j < \pi_i \}\big|
\end{equation}
counts the inversions of $\pi_i$ with respect to the tail. The Lehmer code is
a bijection between permutations and $\{0, 1, \dots, 52! - 1\}$; the resulting
integer is then passed through the extractor.

% ============================================================
\section{Statistical Test Battery}
% ============================================================

Given a byte string $D$ of length $n$, define its bit expansion
$\mathbf{b} = (b_1, \dots, b_{8n}) \in \{0,1\}^{8n}$ via
$b_{8(j-1)+k} = \lfloor d_j / 2^{8-k} \rfloor \bmod 2$.

\subsection{Frequency (monobit) test}

Let $S = \sum_{i=1}^{m} (2 b_i - 1)$ with $m = 8n$. Under the null hypothesis
$b_i \sim \mathrm{Bernoulli}(1/2)$ independently, $S / \sqrt{m}$ converges to a
standard normal. The test statistic is
\begin{equation}
    s_{\text{obs}} \;=\; \frac{|S|}{\sqrt{m}},
    \qquad
    p_{\text{mono}} \;=\; \mathrm{erfc}\!\left( \frac{s_{\text{obs}}}{\sqrt{2}} \right),
    \label{eq:monobit}
\end{equation}
where $\mathrm{erfc}(x) = \frac{2}{\sqrt{\pi}} \int_x^\infty e^{-t^2}\, dt$.
The test passes at significance level $\alpha = 0.01$ when $p_{\text{mono}} \ge \alpha$.

\subsection{Runs test}

Let $V_{\text{obs}} = 1 + \sum_{i=1}^{m-1} \mathbf{1}[b_i \ne b_{i+1}]$ be the
observed number of runs. Under $H_0$, with $\pi = m^{-1} \sum_i b_i$, the
statistic
\begin{equation}
    Z \;=\; \frac{V_{\text{obs}} - 2 m \pi (1 - \pi)}{2 \sqrt{2m}\,\pi(1-\pi)}
    \label{eq:runs}
\end{equation}
is asymptotically standard normal, so
\begin{equation}
    p_{\text{runs}} \;=\; \mathrm{erfc}\!\left( \frac{|V_{\text{obs}}
    - 2 m \pi (1 - \pi)|}{2 \sqrt{2m}\, \pi (1-\pi)} \right).
\end{equation}
A preliminary consistency check $|\pi - 1/2| < 2 / \sqrt{m}$ is required.

\subsection{Chi-squared test on byte frequencies}

Let $c_v$ be the observed count of byte value $v \in \{0, \dots, 255\}$, and
$E_v = n/256$ its expected count. The Pearson statistic is
\begin{equation}
    \chi^2 \;=\; \sum_{v=0}^{255} \frac{(c_v - E_v)^2}{E_v},
    \qquad
    p_{\chi^2} \;=\; 1 - F_{\chi^2_{255}}(\chi^2),
    \label{eq:chi2}
\end{equation}
where $F_{\chi^2_{255}}$ is the CDF of a chi-squared distribution with $255$
degrees of freedom.

\subsection{Autocorrelation}

For lag $\ell$, define
\begin{equation}
    \hat{\rho}(\ell) \;=\;
    \frac{\sum_{i=1}^{n-\ell} (d_i - \bar{d})(d_{i+\ell} - \bar{d})}
         {\sum_{i=1}^{n} (d_i - \bar{d})^2},
    \label{eq:autocorr}
\end{equation}
with $\bar{d} = n^{-1} \sum_i d_i$. For an i.i.d.\ uniform source,
$\hat{\rho}(\ell) \to 0$ in probability as $n \to \infty$.

\subsection{NIST SP 800-22 subset}

The following tests are implemented.

\paragraph{Block frequency.}
Partition $\mathbf{b}$ into $M = \lfloor m / L \rfloor$ blocks of length $L$
(here $L = 128$). Let $\pi_i$ be the proportion of ones in block $i$. Under
$H_0$,
\begin{equation}
    \chi^2 \;=\; 4 L \sum_{i=1}^{M} \left(\pi_i - \tfrac{1}{2}\right)^2
    \;\sim\; \chi^2_M,
    \qquad
    p \;=\; 1 - F_{\chi^2_M}(\chi^2).
    \label{eq:blockfreq}
\end{equation}

\paragraph{Cumulative sums.}
Let $X_i = 2 b_i - 1 \in \{-1, +1\}$ and define the partial sums
$S_k = \sum_{i=1}^{k} X_i$. Let
\begin{equation}
    z \;=\; \max_{1 \le k \le m} |S_k|.
\end{equation}
The $p$-value is computed via the standard normal CDF $\Phi$:
\begin{equation}
    p \;=\; 1 - \sum_{k} \left[ \Phi\!\left(\frac{(4k+1)z}{\sqrt{m}}\right)
    - \Phi\!\left(\frac{(4k-1)z}{\sqrt{m}}\right) \right]
    + \sum_{k} \left[ \Phi\!\left(\frac{(4k+3)z}{\sqrt{m}}\right)
    - \Phi\!\left(\frac{(4k+1)z}{\sqrt{m}}\right) \right],
    \label{eq:cumsu}
\end{equation}
summing over $k$ in the appropriate finite ranges derived in NIST~SP~800-22
\S2.13. The backward mode applies the same formula to the reversed sequence.

\paragraph{Longest run of ones.}
Partition $\mathbf{b}$ into $N = \lfloor m / M \rfloor$ blocks of length $M$
(here $M = 8$). Let $v_j$ count blocks whose longest run of ones belongs to
category $j \in \{0,1,2,3\}$ (max run $\le 1$, $=2$, $=3$, $\ge 4$). Under
$H_0$, the expected proportions are
$(\pi_0, \pi_1, \pi_2, \pi_3) = (0.2148, 0.3672, 0.2305, 0.1875)$
(interpolated for $M = 8$; exact values given in NIST SP 800-22, \S2.4).
The test statistic is
\begin{equation}
    \chi^2 \;=\; \sum_{j=0}^{3} \frac{(v_j - N \pi_j)^2}{N \pi_j}
    \;\sim\; \chi^2_3.
\end{equation}

% ============================================================
\section{Composite Scoring Functional}
% ============================================================

Let $M_i$ be a method and let $\mathbf{f}(M_i) = (f_1, \dots, f_7)$ be its
aggregated metrics:
\begin{equation}
\begin{aligned}
    f_1 &= \frac{\hat{H}_\infty(D)}{8}, &\quad
    f_2 &= \frac{\hat{H}_2(D)}{8}, \\
    f_3 &= \mathbf{1}[p_{\text{mono}} \ge \alpha], &\quad
    f_4 &= \mathbf{1}[p_{\text{runs}} \ge \alpha], \\
    f_5 &= \mathbf{1}[p_{\chi^2} \ge \alpha], &\quad
    f_6 &= 1 - \min(1, |\hat{\rho}(1)|), \\
    f_7 &= \rho(D).
\end{aligned}
\label{eq:features}
\end{equation}
Each $f_j \in [0,1]$. The weighted score is
\begin{equation}
    S_0(M_i) \;=\; 100 \cdot \sum_{j=1}^{7} w_j\, f_j(M_i),
    \qquad
    \sum_j w_j = 1,
    \label{eq:score0}
\end{equation}
with weights
\begin{equation}
    (w_1, \dots, w_7) = (0.30,\, 0.10,\, 0.15,\, 0.15,\, 0.10,\, 0.10,\, 0.10),
\end{equation}
corresponding to (min-entropy, collision, monobit, runs, chi-squared,
autocorrelation, compression).

\subsection{Penalization for determinism}

\begin{definition}[Determinism penalty]
Let $\delta(M_i) \in \{0,1\}$ indicate that $M_i$ is deterministic and let
$\beta(M_i) \in \{0,1\}$ indicate that its theoretical entropy per unit is zero.
The final score is
\begin{equation}
    S(M_i) \;=\; S_0(M_i) \cdot
    \begin{cases}
        0.5 & \text{if } \delta(M_i) = 1, \\
        1.0 & \text{otherwise},
    \end{cases}
    \;\cdot\;
    \begin{cases}
        0.75 & \text{if } \beta(M_i) = 1, \\
        1.0  & \text{otherwise}.
    \end{cases}
    \label{eq:score}
\end{equation}
\end{definition}

\begin{remark}
The penalties encode a cryptographic axiom: a method that merely \emph{looks}
random is not equivalent to a method that \emph{is} unpredictable. A
deterministic chaotic map with a known seed has $H_\infty = 0$ regardless of how
well its output passes the statistical battery.
\end{remark}

% ============================================================
\section{Cryptographic Pipeline}
% ============================================================

\subsection{Seed derivation}

Given raw entropy $E \in \{0,1\}^{8n}$ and an optional passphrase $P$, the seed
is derived via PBKDF2 with HMAC-SHA512:
\begin{equation}
    S \;=\; \mathrm{PBKDF2\text{-}HMAC\text{-}SHA512}
    \big(E,\; \texttt{"mnemonic"} \| P,\; c = 2048,\; \mathrm{dkLen} = 64\big).
    \label{eq:pbkdf2}
\end{equation}
$S$ is a 512-bit seed suitable for BIP32 key derivation.

\subsection{BIP32 master key}

Let $\mathrm{HMAC}(K, \cdot)$ denote HMAC-SHA512 with key $K$. The master key
pair is
\begin{equation}
    I \;=\; \mathrm{HMAC}\big(\texttt{"Bitcoin seed"},\; S\big),
    \qquad
    k_m \;=\; I_{[0{:}32]}, \qquad
    c_m \;=\; I_{[32{:}64]},
    \label{eq:bip32master}
\end{equation}
where $k_m$ is the master private key and $c_m$ the master chain code.

\subsection{Hardened child derivation}

For a hardened child with index $i \in [0, 2^{31})$:
\begin{equation}
    I \;=\; \mathrm{HMAC}\big(c_m,\; \texttt{0x00} \| k_m \| i_{\text{BE32}}\big),
    \qquad
    k_i \;=\; I_{[0{:}32]}, \qquad
    c_i \;=\; I_{[32{:}64]}.
    \label{eq:bip32child}
\end{equation}
The hardened mode prevents deriving child public keys without the parent
private key, which is required for cold-storage signers.

\subsection{Multi-signature script}

Given three signer public keys $\mathrm{PK}_1, \mathrm{PK}_2, \mathrm{PK}_3$
(determined by their respective $k_i$ on secp256k1), a $2$-of-$3$ P2SH
multisignature script is
\begin{equation}
    \texttt{Script} \;=\; \mathrm{OP}_2 \;
    \langle \mathrm{PK}_1 \rangle \langle \mathrm{PK}_2 \rangle
    \langle \mathrm{PK}_3 \rangle \; \mathrm{OP}_3 \; \mathrm{OP\_CHECKMULTISIG}.
    \label{eq:multisig}
\end{equation}
A transaction spending an output locked by this script requires $m = 2$
distinct valid signatures. Formally, if $\sigma_i$ is a valid ECDSA signature
under $\mathrm{PK}_i$ for message $T$, the script accepts if and only if
\begin{equation}
    \exists\, \mathcal{S} \subseteq \{1,2,3\},\; |\mathcal{S}| \ge 2,\;
    \forall i \in \mathcal{S}:\; \mathrm{Verify}(\mathrm{PK}_i, T, \sigma_i) = 1.
\end{equation}

\subsection{Security of independent sources}

Suppose the multisignature uses $k$ independent entropy sources with
min-entropies $H_\infty^{(1)}, \dots, H_\infty^{(k)}$. If each source is
adversarially controlled except one with min-entropy $H_\infty^* \ge 128$ bits,
the resulting wallet remains $128$-bit secure. Formally, define the guessing
probability of the adversary as
\begin{equation}
    P_{\text{guess}} \;=\; \max_{\mathcal{A}} \Pr[\mathcal{A} \text{ recovers key}].
\end{equation}
Then
\begin{equation}
    P_{\text{guess}} \;\le\; 2^{-H_\infty^*} \cdot
    \prod_{j \ne *} 1 \;=\; 2^{-H_\infty^*}.
    \label{eq:multisig_security}
\end{equation}
This is the formal justification for the multisig design: even if $k-1$ sources
are compromised, the wallet's security reduces to that of the single honest
source.

% ============================================================
\section{Discussion}
% ============================================================

The framework presented here has a deliberate pedagogical asymmetry. On one
hand, every method --- whether a CSPRNG, a chaotic map, or a physical die --- is
subjected to the same statistical battery, and many will pass all tests. On the
other hand, the scoring functional strictly separates methods that
\emph{contribute} entropy from those that merely \emph{propagate} it. This
separation is captured by the determinism penalty of
Equation~\eqref{eq:score} and by the use of min-entropy as the dominant term in
Equation~\eqref{eq:score0}.

Three formal consequences follow:
\begin{enumerate}[leftmargin=1.4em]
    \item \textbf{Statistical uniformity is necessary but not sufficient.}
          Any deterministic system with a secret seed can be made statistically
          indistinguishable from a true RNG; this is the entire purpose of a
          hash function used as an extractor.
    \item \textbf{Physical entropy scales linearly with effort.} Generating
          $128$ bits from dice requires at least $\lceil 128 / \log_2 6 \rceil
          = 50$ rolls; from coin flips, $128$ tosses; from a deck shuffle, at
          most one shuffle suffices.
    \item \textbf{Multisig security is governed by the minimum, not the sum.}
          Equation~\eqref{eq:multisig_security} formalizes the well-known
          principle that in threshold cryptography, the adversary only needs to
          break one honest participant.
\end{enumerate}

The arena thus serves two purposes: as an empirical benchmark for individual
generators, and as a formal demonstration that the cryptographic quality of a
key depends on a property --- unpredictability under adversarial observation
--- that no statistical test alone can measure.

% ============================================================
\appendix
\section{Notation}
% ============================================================

\begin{table}[h]
\centering
\begin{tabular}{@{}ll@{}}
\toprule
Symbol & Meaning \\
\midrule
$X, Y$ & Discrete random variables \\
$p(x)$ & Probability mass function of $X$ \\
$H_1, H_2, H_\infty$ & Shannon, collision, and min-entropy \\
$\hat{H}$ & Empirical estimate of $H$ \\
$D$ & Byte string output by a generator \\
$n, m$ & Number of bytes, number of bits ($m = 8n$) \\
$r$ & Logistic map control parameter \\
$\sigma, \rho, \beta$ & Lorenz parameters \\
$\lambda$ & Lyapunov exponent \\
$\alpha$ & Significance level of statistical tests \\
$w_j$ & Weight of the $j$-th metric in the score \\
$S(M_i)$ & Final score of method $M_i$ \\
$k_m, c_m$ & BIP32 master private key and chain code \\
$\mathrm{PK}_i$ & Public key of signer $i$ \\
\bottomrule
\end{tabular}
\end{table}

\section{Constants}
% ============================================================

\begin{table}[h]
\centering
\begin{tabular}{@{}lr@{}}
\toprule
Quantity & Value \\
\midrule
$\log_2 6$ & $2.58496\ldots$ \\
$\log_2(52!)$ & $225.58100\ldots$ \\
$52!$ & $\approx 8.0658 \times 10^{67}$ \\
MT19937 period & $2^{19937} - 1$ \\
secp256k1 group order & $\approx 1.1579 \times 10^{77}$ \\
Lorenz largest Lyapunov exponent & $\approx 0.9056$ \\
Logistic map Lyapunov at $r=4$ & $\ln 2$ \\
PBKDF2 iterations (BIP39) & $2048$ \\
\bottomrule
\end{tabular}
\end{table}

\end{document}